\section{Experimental Simulation Results and Analytical Bounds}

Having executed the Object-Oriented emulation framework across discrete-event simulation sweeps, we generate empirical performance data evaluating the proposed Frugal Edge architecture against legacy baseline paradigms under progressive Electronic Warfare (EW) and RF degradation. We analyze the system across three core performance dimensions: Effective State Synchronization, Communication Bandwidth Efficiency, and Partition Resilience.

\subsection{Effective State Synchronization Under Progressive EW Degradation}
We evaluate three distinct protocol architectures across nine progressive network degradation levels, ranging from benign conditions ($p=0.0$) to catastrophic RF packet loss ($p=0.80$):
\begin{enumerate}
    \item \textbf{Gossip Protocol (Baseline 1):} Blind $1$-hop broadcast of uncompressed state matrices ($\approx 450$ bytes per payload) with a maximum hop limit of $\text{TTL}=3$.
    \item \textbf{Epidemic Routing Protocol (Baseline 2):} Store-and-forward flooding to all unvisited encountered neighbors without TTL constraints, maximizing delivery redundancy at the cost of high bandwidth consumption.
    \item \textbf{Agentic SLM Protocol (Proposed):} Quantized semantic payload compression ($\approx 32\text{--}48$ bytes) coupled with Exponential Moving Average (EMA) link-reliability scoring and dynamic 2-hop relay selection.
\end{enumerate}

\begin{table}[h]
\centering
\caption{Empirical State Synchronization (\%) and Delivery Rate (\%) Across RF Drop Rates}
\label{tab:empirical_sync}
\resizebox{\linewidth}{!}{%
\begin{tabular}{@{}cffffff@{}}
\toprule
& \multicolumn{2}{c}{\textbf{Gossip (Baseline 1)}} & \multicolumn{2}{c}{\textbf{Epidemic (Baseline 2)}} & \multicolumn{2}{c}{\textbf{Agentic SLM (Proposed)}} \\
\cmidrule(lr){2-3} \cmidrule(lr){4-5} \cmidrule(lr){6-7}
\textbf{Drop Rate ($p$)} & \textbf{Sync (\%)} & \textbf{Delivery (\%)} & \textbf{Sync (\%)} & \textbf{Delivery (\%)} & \textbf{Sync (\%)} & \textbf{Delivery (\%)} \\ \midrule
0\% (Benign)  & 96.6\% & 99.7\% & 97.4\% & 96.9\% & 96.6\% & \mathbf{99.7\%} \\
10\%          & 94.6\% & 88.6\% & 97.1\% & 88.5\% & 96.9\% & \mathbf{97.2\%} \\
20\%          & 94.7\% & 80.0\% & 96.6\% & 77.4\% & 95.5\% & \mathbf{94.2\%} \\
30\%          & 86.1\% & 68.1\% & 96.7\% & 69.7\% & 96.3\% & \mathbf{93.8\%} \\
40\%          & 79.6\% & 59.0\% & 95.6\% & 57.7\% & 94.8\% & \mathbf{90.0\%} \\
50\%          & 67.4\% & 49.0\% & 95.9\% & 48.3\% & 93.8\% & \mathbf{88.6\%} \\
60\%          & 54.2\% & 39.2\% & 94.4\% & 39.8\% & 93.1\% & \mathbf{85.8\%} \\
70\%          & 40.5\% & 29.5\% & 93.7\% & 28.4\% & 91.2\% & \mathbf{84.7\%} \\
80\% (Severe) & 27.0\% & 18.0\% & 90.9\% & 19.5\% & \mathbf{91.6\%} & \mathbf{83.2\%} \\ \bottomrule
\end{tabular}%
}
\end{table}

As detailed in Table \ref{tab:empirical_sync}, the Gossip Protocol experiences rapid performance degradation as RF loss escalates, dropping from $96.6\%$ state synchronization at $p=0.0$ to a failure threshold of $27.0\%$ at $p=0.80$. Epidemic Routing maintains high synchronization ($90.9\%$ at $p=0.80$) due to brute-force packet duplication, but suffers low direct delivery success ($19.5\%$). 

In contrast, the proposed Agentic SLM Protocol maintains $91.6\%$ global state synchronization and an $83.2\%$ delivery success rate under severe $80\%$ packet loss. By shrinking payload sizes by an $8\times$ factor and dynamically routing around degraded edges, the proposed protocol prevents the delivery failure characteristic of uncompressed baseline broadcasts.

\subsection{Communication Overhead and Byte Efficiency}
A central requirement for edge deployment in D-DIL environments is minimizing physical layer throughput. Table \ref{tab:empirical_bytes} summarizes the aggregate network traffic generated by each protocol across the simulation duration.

\begin{table}[h]
\centering
\caption{Aggregate Transmitted Bytes Across Protocol Architectures}
\label{tab:empirical_bytes}
\begin{tabular}{@{}crrr@{}}
\toprule
\textbf{Drop Rate ($p$)} & \textbf{Gossip (B)} & \textbf{Epidemic (B)} & \textbf{Agentic SLM (B)} \\ \midrule
0\% (Benign)  & 1,751,668 & 45,358,865 & 642,733 \\
20\%          & 1,338,994 & 35,457,627 & 594,876 \\
40\%          &   815,119 & 25,770,564 & 551,843 \\
60\%          &   399,167 & 17,088,884 & 512,510 \\
80\% (Severe) &   108,473 & 7,630,730  & \mathbf{474,599} \\ \bottomrule
\end{tabular}
\end{table}

At an $80\%$ packet loss rate, Epidemic Routing consumes $7,630,730$ bytes ($\approx 7.63$ MB) of network traffic to sustain state synchronization. In comparison, the Agentic SLM Protocol consumes only $474,599$ bytes ($\approx 474$ KB). The proposed architecture achieves comparable state synchronization to Epidemic Routing ($91.6\%$ vs. $90.9\%$) while achieving a \textbf{$16.1\times$ reduction in total network byte overhead}.

\subsection{Scenario Walkthrough: Partition Resilience}
Consider an operational scenario where physical obstacles or localized EW splitting partition a $20$-node swarm into two disjoint sub-graphs of $10$ nodes each. In a centralized system, any sub-graph lacking a link to the central server becomes inert. 

Under the proposed decentralized auction protocol, the adjacency matrix $A(t)$ becomes block-diagonal. The algorithm guarantees that both sub-swarms independently execute consensus auctions within their respective reachable topologies. When communication links re-establish and $A(t)$ becomes connected, the winning bid vectors ($\text{WBL}$) reconcile asynchronously, restoring unified global consensus in $O(D)$ time without requiring state resets.

\section{Discussion and Conclusion}

The empirical findings demonstrate that operating high-bandwidth or centralized telemetry architectures in D-DIL environments introduces severe operational vulnerabilities. While previous research has focused on federated model training on edge devices \cite{imteaj2022}, real-time operational inference requires low-latency, deterministic allocation protocols.

By incorporating localized agentic reasoning engines and extreme payload quantization \cite{steinberger2026}, the Frugal Edge architecture eliminates single points of failure. The empirical benchmark confirms that under an $80\%$ RF packet drop rate, the Agentic SLM Protocol achieves $91.6\%$ state synchronization while reducing network byte consumption by $16.1\times$ relative to Epidemic Routing paradigms. Future research will explore scaling laws across non-stationary Gauss-Markov mobility topologies and hardware-in-the-loop transceiver validation.

\begin{thebibliography}{99}

\bibitem{bekmezci2013}
I. Bekmezci, O. K. Sahingoz, and S. Temel, ``Flying ad-hoc networks (FANETs): A survey,'' \textit{Ad Hoc Networks}, vol. 11, no. 3, pp. 1254--1270, 2013.

\bibitem{choi2009}
H. L. Choi, L. Brunet, and J. P. How, ``Consensus-based decentralized auctions for robust task allocation,'' \textit{IEEE Transactions on Robotics}, vol. 25, no. 4, pp. 912--926, 2009.

\bibitem{imteaj2022}
A. Imteaj, U. Thakker, S. Wang, J. Li, and M. H. Amini, ``A survey on federated learning for resource-constrained IoT devices,'' \textit{IEEE Internet of Things Journal}, vol. 9, no. 1, pp. 4756--4789, 2022.

\bibitem{junior2025}
N. F. Junior, \textit{et al.}, ``FedSensor: Federated learning framework for secure sensor-based IoT at the extreme edge,'' \textit{IEEE Access}, vol. 13, pp. 136940--136955, 2025.

\bibitem{kairouz2021}
P. Kairouz, \textit{et al.}, ``Advances and open problems in federated learning,'' \textit{Foundations and Trends in Machine Learning}, vol. 14, no. 1-2, pp. 1--210, 2021.

\bibitem{mcmahan2017}
H. B. McMahan, E. Moore, D. Ramage, S. Hampson, and B. A. y Arcas, ``Communication-efficient learning of deep networks from decentralized data,'' \textit{Proceedings of the 20th International Conference on Artificial Intelligence and Statistics (AISTATS)}, vol. 54, pp. 1273--1282, 2017.

\bibitem{shi2016}
W. Shi, J. Cao, Q. Zhang, Y. Li, and L. Xu, ``Edge computing: Vision and challenges,'' \textit{IEEE Internet of Things Journal}, vol. 3, no. 5, pp. 637--646, 2016.

\bibitem{steinberger2026}
P. Steinberger, ``OpenClaw — Personal AI assistant,'' GitHub Repository, 2026. [Online]. Available: \url{https://github.com/openclaw/openclaw}

\bibitem{suri2016}
N. Suri, M. Tortonesi, J. Michaelis, P. Budulas, G. Benincasa, S. Russell, C. Stefanell, and R. Winkler, ``Analyzing the applicability of Internet of Things to the battlefield environment,'' in \textit{2016 International Conference on Military Communications and Information Systems (ICMCIS)}, 2016, pp. 1--8.

\end{thebibliography}

\end{document}
